\chapter{2010:Richard F. Heck and Ei-ichi Negishi and Akira Suzuki}

\label{chap:chemistry2010}

The Nobel Prize in Chemistry for the year 2010 was awarded jointly to Richard F. Heck and Ei-ichi Negishi and Akira Suzuki.

\textbf{Motivation:}

for palladium-catalyzed cross couplings in organic synthesis

\section{Richard F. Heck}

\subsection*{Early Life and Education}

Richard F. Heck was born in 1931 in Springfield, Massachusetts, USA.

\subsection*{Career and Major Research}

Heck received his doctorate from the University of California, Los Angeles, in 1954. He worked for the Hercules Company and later became a professor at the University of Delaware, where he developed the palladium-catalyzed coupling of organic halides with alkenes now known as the Heck reaction.

\subsection*{Nobel Prize-winning Work}

The work for which Richard F. Heck was awarded the Nobel Prize in Chemistry in 2010 was described by the Nobel Committee as: "for palladium-catalyzed cross couplings in organic synthesis".

The Heck reaction joins an aryl or vinyl halide to an alkene in the presence of a palladium catalyst, forming a new carbon--carbon bond.

\subsection*{Significance and Impact}

The Heck reaction was the first practical way to form carbon--carbon bonds between complex molecules under mild conditions. It is used throughout the pharmaceutical and materials industries.

\subsection*{Later Career and Honors}

Heck retired from the University of Delaware and later moved to the Philippines. He died on 2015-10-09 in Manila, Philippines.

\subsection*{Legacy}

The Heck reaction remains one of the most widely used transformations in organic synthesis and is the basis of many industrial processes.

\subsection*{Selected Publications}

"Palladium-Catalyzed Reactions of Organic Halides with Olefins" (Accounts of Chemical Research, 1979)

\clearpage

\section{Ei-ichi Negishi}

\subsection*{Early Life and Education}

Ei-ichi Negishi was born in 1935 in Changchun, China.

\subsection*{Career and Major Research}

Negishi earned a bachelor's degree from the University of Tokyo in 1958 and a doctorate from the University of Pennsylvania in 1963. He worked as a postdoctoral researcher with Herbert C. Brown at Purdue University and later became a professor there. In the late 1970s he developed the palladium- or nickel-catalyzed cross coupling of organozinc compounds with organic halides.

\subsection*{Nobel Prize-winning Work}

The work for which Ei-ichi Negishi was awarded the Nobel Prize in Chemistry in 2010 was described by the Nobel Committee as: "for palladium-catalyzed cross couplings in organic synthesis".

The Negishi coupling uses an organozinc reagent to transfer a carbon group to a palladium or nickel catalyst, which then delivers it to an organic halide.

\subsection*{Significance and Impact}

The Negishi coupling allows carbon--carbon bonds to be made with high selectivity and is widely used in the synthesis of pharmaceuticals, agrochemicals, and materials.

\subsection*{Later Career and Honors}

Negishi remained a professor at Purdue University for four decades. He died on 2021-06-06 in Indianapolis, Indiana, USA.

\subsection*{Legacy}

The Negishi coupling is a standard reaction of modern cross-coupling chemistry and is taught and used in laboratories worldwide.

\subsection*{Selected Publications}

"Palladium- or Nickel-Catalyzed Cross Coupling: A New Selective Method for Carbon-Carbon Bond Formation" (Accounts of Chemical Research, 1982)

\clearpage

\section{Akira Suzuki}

\subsection*{Early Life and Education}

Akira Suzuki was born in 1930 in Mukawa, Japan.

\subsection*{Career and Major Research}

Suzuki received his doctorate from Hokkaido University in 1959 and became a professor there. He developed the cross coupling of organoboron compounds, such as boronic acids, with organic halides in the presence of a palladium catalyst, a reaction now known as the Suzuki coupling.

\subsection*{Nobel Prize-winning Work}

The work for which Akira Suzuki was awarded the Nobel Prize in Chemistry in 2010 was described by the Nobel Committee as: "for palladium-catalyzed cross couplings in organic synthesis".

The Suzuki coupling forms carbon--carbon bonds between an organoboron compound and an organic halide, with the boron group transferred under mild conditions.

\subsection*{Significance and Impact}

Because organoboron reagents are stable and easy to handle, the Suzuki coupling is one of the most widely used reactions for building complex molecules, including pharmaceuticals and electronic materials.

\subsection*{Later Career and Honors}

Suzuki is professor emeritus of Hokkaido University and continues to be active in chemistry.

\subsection*{Legacy}

The Suzuki coupling is a cornerstone of cross-coupling chemistry and is used in industrial processes ranging from drug manufacture to the production of liquid crystal displays.

\subsection*{Selected Publications}

"Stereoselective Synthesis of Arylated (E)-Alkenes by the Reaction of Alk-1-enylboranes with Aryl Halides in the Presence of Palladium Catalyst" (Journal of the Chemical Society, Chemical Communications, 1979)

\clearpage

\section*{Chapter Summary}

The 2010 prize recognized the palladium-catalyzed cross-coupling reactions developed by Heck, Negishi, and Suzuki, which make it possible to join carbon atoms selectively. These reactions are now indispensable in the production of pharmaceuticals, agrochemicals, and advanced materials.
